\documentclass[12pt]{article}

% ---------- Encoding & typography ----------
\usepackage[utf8]{inputenc}
\usepackage[T1]{fontenc}
\usepackage{lmodern}
\usepackage{microtype}
\usepackage{natbib}

% ---------- Layout ----------
\usepackage[margin=1in]{geometry}
\usepackage{setspace}
\onehalfspacing

% ---------- Math & symbols ----------
\usepackage{amsmath,amssymb}

% ---------- Headers ----------
\usepackage{fancyhdr}
\setlength{\headheight}{14pt}
\pagestyle{fancy}
\fancyhf{}
\fancyhead[L]{The Entitlement Cascade}
\fancyhead[R]{J.\ Christian McKinley}
\fancyfoot[C]{\thepage}

% ---------- Links ----------
\usepackage{hyperref}
\hypersetup{
  colorlinks=true,
  linkcolor=blue,
  urlcolor=blue,
  citecolor=blue
}

\title{\textbf{The Entitlement Cascade}\\
\large Why Future Theft Cannot Complete Past Restitution}
\author{J.\ Christian McKinley}
\date{12 May 2026}

\begin{document}
\maketitle

\begingroup
  \footnotetext[0]{draft}
\endgroup

\begin{abstract}
This paper states a narrow cascade-side no-go: a coercive transfer chain cannot become moral by using each new taking to satisfy the prior victim class. In a serial structure where Peter is wronged, Paul is later coerced to compensate Peter, Joseph is later coerced to compensate Paul, and so on, each apparent act of repair creates the next grievance. The result is not restitution but entitlement cascade: a sequence in which every prior extraction becomes the rhetorical basis for the next extraction. The system purchases peace with one claimant class by creating another claimant class downstream. The familiar pay-as-you-go defense follows this pattern: one cohort is told it is owed because it was previously forced to pay, and the next cohort is forced to satisfy that claim. Continuing the chain cannot complete the injustice because every continuation generates the next unpaid injury. Full unwind may be theoretically clean but practically unavailable in mass systems. The only non-propagating remedy is cessation: stop the chain, recognize the final aggrieved class, and refuse to create new victims to hide the prior wrong.
\end{abstract}

\section{Introduction}

Future theft cannot complete past restitution.

A coercive transfer chain creates a special moral problem. The first taking creates a victim. The next taking is justified as compensation for the first victim. That second taking creates a new victim. The third taking is justified as compensation for the second victim. The chain continues, with each round of extraction producing the next entitlement claim.

This is not restitution. It is cascade.

The structure can be stated simply:

\begin{quote}
Steal from Peter. Later steal from Paul to pay Peter. Later steal from Joseph to pay Paul. Later steal from another to pay Joseph.
\end{quote}

Each payment appears to answer the prior grievance. In fact, each payment creates the next grievance. The system never repairs the wrong from clean resources. It shifts the wrong forward.

This paper calls that structure the entitlement cascade.

The central claim is narrow:

\begin{quote}
A coercive transfer chain cannot be morally completed by continuation, because every continuation creates the next claimant class.
\end{quote}

The only non-propagating remedy is cessation. Stop the chain. Name the final uncompensated injury. Refuse to create a new victim in order to disguise it.

\section{The Basic Chain}

Consider the simplest case.

Agent \(A\) steals from Peter. Peter is wronged. Later, \(A\) steals from Paul and gives the proceeds to Peter. Peter now receives compensation in appearance, but Paul has been wronged. Later, \(A\) steals from Joseph and gives the proceeds to Paul. Paul now receives compensation in appearance, but Joseph has been wronged.

The pattern is:

\[
P_1 \leftarrow P_2 \leftarrow P_3 \leftarrow P_4 \leftarrow \cdots
\]

where each \(P_n\) is paid by extraction from \(P_{n+1}\).

At no point does the wrongdoer repair the wrong from clean resources. At each point, one grievance is quieted by creating another. The debt is not paid. It is rolled forward into the next victim.

\section{Entitlement Cascade Defined}

\textbf{Definition 1.} An entitlement cascade is a serial coercive transfer structure in which each recipient's claim is justified by prior extraction from that recipient, while the funds used to satisfy the claim are extracted from a later or current party who thereby becomes the next claimant.

The key features are:

\begin{enumerate}
  \item prior extraction creates a grievance;
  \item the grievance is redescribed as entitlement;
  \item the entitlement is funded by new extraction;
  \item the new extraction creates another grievance;
  \item the new grievance becomes the next entitlement.
\end{enumerate}

This structure does not end through moral completion. It ends only by collapse, reform, refusal, or open acknowledgment of a final unpaid injury.

\section{Axioms}

The argument rests on eight axioms.

\begin{enumerate}
  \item \textbf{Wrongful extraction creates a claim against the extractor.} The victim of a taking may have a claim against the party that took.

  \item \textbf{A claim against the extractor does not automatically transfer to later parties.} The fact that Peter was wronged by \(A\) does not make Paul Peter's debtor.

  \item \textbf{Restitution requires clean source.} Compensation funded by new wrongful taking is not clean restitution.

  \item \textbf{Serial taking serializes injury.} Each new extraction creates a new victim or claimant class.

  \item \textbf{Continuation does not cure cascade.} A chain that requires the next victim to compensate the prior victim cannot become moral by adding another link.

  \item \textbf{Legal entitlement does not erase cascade.} A statute declaring each claimant eligible does not make the funding source clean.

  \item \textbf{A final grievance is not a warrant for the next grievance.} The fact that stopping leaves one class uncompensated does not justify creating another class of victims.

  \item \textbf{Cessation is not full repair.} Stopping the chain prevents propagation; it does not pretend that all prior injuries have been repaired.
\end{enumerate}

These axioms distinguish moral completion from mere postponement.

\section{Proposition: The Cascade Cannot Be Completed by Continuation}

\textbf{Proposition.} In a serial coercive transfer chain where each claimant class is paid from funds coercively extracted from a later class, the chain cannot be morally completed by continuation.

\textbf{Proof.} Let \(C_n\) be a claimant class that was previously coerced from, and let \(C_{n+1}\) be the later class coerced to fund payment to \(C_n\). By Axiom 1, \(C_n\) may possess a claim against the extractor. By Axiom 2, that claim does not automatically lie against \(C_{n+1}\). If \(C_{n+1}\) is coerced to pay \(C_n\), then by Axiom 3 the payment is not clean restitution. By Axiom 4, the coercion of \(C_{n+1}\) creates a new claimant class. Therefore the continuation from \(C_n\) to \(C_{n+1}\) does not complete the chain; it extends it.

Because every additional continuation creates another wrong of the same kind, the chain cannot be morally completed by adding more links. \(\square\)

\section{The Musical Chairs Problem}

The cascade is moral musical chairs.

Every round keeps the music playing by taking a chair from the next person. As long as the music continues, someone appears to be seated. But the system does not create enough chairs. It delays the moment when the final person is left standing.

Stopping the music reveals the final injury. Continuing the music does not eliminate that injury. It merely decides that someone else will be left standing later.

This matters because defenders of cascade systems often say:

\begin{quote}
If we stop now, this cohort will be harmed.
\end{quote}

That is true. But it is not a defense of continuation. It is a description of the chain's final unpaid injury. Continuing the chain only creates a later final unpaid injury in another class.

The moral question is not whether stopping is painless. It is whether pain can be eliminated by creating the next victim. It cannot.

\section{The Ponzi Structure}

The entitlement cascade has a structural kinship with a Ponzi scheme.

In a Ponzi scheme, early participants are paid from later participants' money rather than from real returns. The early payments create legitimacy. They prove the scheme appears to work. Later entrants are drawn in or held in place by the expectation that the same structure will eventually favor them.

The public cascade differs in legal status and enforcement mechanism. Participation may be compulsory rather than fraudulently induced. The state may define the system as benefit, insurance, public obligation, or entitlement rather than investment. But the structural dependency remains: earlier claimants are satisfied with later inflows.

The relevant moral point is not whether the public system satisfies every legal element of private fraud. The point is that each payment is funded by the next extraction.

A Ponzi scheme does not become honest because early victims are paid. It becomes larger.

\section{Pay-As-You-Go Transfer}

Pay-as-you-go transfer programs provide the clearest public case.

One cohort is compelled to pay while working. Those funds are used to pay prior beneficiaries. Later, that paying cohort becomes the beneficiary class. It then claims payment because it paid in. But its payment comes from a new cohort of workers currently being compelled to pay.

The defense says:

\begin{quote}
They paid in, so they are owed.
\end{quote}

The cascade analysis responds:

\begin{quote}
They may have a claim against the system that took from them. They do not thereby acquire title against the next cohort being taken from.
\end{quote}

Each cohort is first victim, then beneficiary, then alibi. The system converts prior extraction into a demand for future extraction. That is the entitlement cascade.

\section{The Claim Conversion}

The cascade depends on converting claim direction.

The original direction is:

\begin{quote}
The prior victim has a claim against the extractor.
\end{quote}

The cascade converts it into:

\begin{quote}
The prior victim has a claim against the next extracted class.
\end{quote}

That conversion is the moral error.

Peter's claim against the thief does not become a claim against Paul. Paul's later claim against the thief does not become a claim against Joseph. Each victim has a claim against the party that structured the taking. The chain survives by redirecting claims away from the wrongdoer and toward the next available payer.

\section{The Rhetoric of ``I Paid In''}

The phrase ``I paid in'' powers the entitlement cascade.

The statement is true as history. The recipient was compelled to contribute. He lost money. He relied on future payment. He may have arranged his life around the promise.

But ``I paid in'' does not answer the source question. It does not show that the current payer owes. It does not prove that the funds now received are the recipient's retained contributions. It does not transform new extraction into restitution.

The cleaner statement is:

\begin{quote}
I was forced to pay into a chain that now proposes to pay me by forcing someone else.
\end{quote}

That sentence preserves the structure. It blocks the cascade's moral laundering.

\section{The Rhetoric of ``We Cannot Stop Now''}

The cascade also depends on the phrase:

\begin{quote}
We cannot stop now.
\end{quote}

The reason given is that stopping would harm those who already paid. Again, the premise is true. The conclusion does not follow.

Stopping exposes the unpaid injury of the final class. Continuing creates another unpaid injury in a later class. The argument against stopping is therefore an argument for relocating the final injury, not eliminating it.

The system says:

\begin{quote}
If we stop, this class loses.
\end{quote}

The answer is:

\begin{quote}
If you continue, another class loses later.
\end{quote}

Continuation is not compassion. It is postponement through fresh coercion.

\section{The Three Options}

Once an entitlement cascade exists, there are three conceptual options:

\begin{enumerate}
  \item full unwind;
  \item continuation;
  \item cessation.
\end{enumerate}

Each option has a different moral status.

Full unwind attempts to reverse the chain. Continuation keeps the chain moving. Cessation stops the creation of new victims while leaving the final injury visible.

Only cessation prevents propagation where full unwind is unavailable.

\section{Option One: Full Unwind}

Full unwind is theoretically clean in small cases.

If a thief stole from Peter, then stole from Paul to pay Peter, then stole from Joseph to pay Paul, a complete court-supervised unwind could attempt to trace funds, reverse transfers, restore title, and impose liability on the thief. In small chains with identifiable parties and traceable assets, unwind may approximate justice.

Mass public systems are different.

Funds have been spent. Recipients may be dead. Payers are dispersed. Benefits are consumed. Records may be incomplete. Reliance has formed. Generations have passed. The wrongdoer may be the state itself, which has no clean treasury separate from those it taxes.

Full unwind may describe an ideal of title correction. It is usually unavailable as a practical mass remedy.

\section{Option Two: Continuation}

Continuation keeps paying the prior class by extracting from the next class.

This is politically attractive because it avoids immediate confrontation. The current beneficiaries are paid. The current payers are told they will be future beneficiaries. The system preserves peace by promising everyone a later seat.

But continuation does not repair. It propagates.

Each new payment creates a new class that can later say:

\begin{quote}
I paid in; now I am owed.
\end{quote}

The chain feeds on this sentence. Every new extraction becomes the moral seed of the next entitlement claim.

Continuation is therefore not a solution. It is the mechanism by which the problem survives.

\section{Option Three: Cessation}

Cessation stops the chain.

It says: no new class will be forced to fund the prior class merely because the prior class was itself forced earlier. The final aggrieved class is recognized as injured. Its injury is not denied. But its injury is not used as authority to create the next victim class.

Cessation is morally severe because it leaves loss visible. It refuses the comforting illusion that everyone can be made whole by moving the cost forward.

The principle is:

\begin{quote}
Stopping does not make the final injury just. It prevents that injury from being used as a warrant to create the next one.
\end{quote}

That is the non-propagating solution.

\section{The Final Aggrieved Class}

Every cascade has a final aggrieved class if it stops.

This class paid into the chain and does not receive the promised later extraction from others. Its grievance is real. It may be angry. It may say truthfully that it was wronged.

But this fact does not justify creating another victim class.

The final aggrieved class is the exposed cost of stopping a chain that should not have begun. The injustice is not created by cessation. It is revealed by cessation. It was already embedded in the structure, waiting for the music to stop.

A system that can stop only by exposing an unpaid victim was never clean.

\section{The Wrong of Hiding the Final Victim}

Continuation hides the final victim by postponing him.

The person left standing is not eliminated. He is moved into the future. The moral burden is transferred to people who may be younger, unborn, politically weaker, or not yet aware of the debt being assigned to them.

This is one reason cascade systems persist. The current beneficiary is visible. The future victim is abstract. The visible claim defeats the invisible cost.

Cessation forces the final victim into view. That is painful. It is also honest.

\section{The Intergenerational Alibi}

The cascade is strongest when organized across generations.

Each generation is told:

\begin{quote}
Pay now; receive later.
\end{quote}

The promise appears reciprocal. But the reciprocity is not clean because participation is compulsory and funding comes from later coerced contributors. Each generation's receipt depends on repeating the wrong against the next.

This creates an intergenerational alibi. The current recipient says he paid. The current payer is told he will be paid. The system hides the fact that the promise requires future coercion.

A voluntary intergenerational contract could be clean. A coerced intergenerational chain is not.

\section{The Cascade and Legal Entitlement}

Legal entitlement stabilizes the cascade.

Each class is given statutory eligibility. The law declares who qualifies, how much is owed, and when payment occurs. This legal status makes the claim appear complete.

But legal entitlement does not solve the source problem. It only formalizes the cascade. It tells Peter that Paul must pay and later tells Paul that Joseph must pay. The statutory form repeats the assignment without proving moral title.

Law can keep the music playing. It cannot make the chairs multiply.

\section{The Cascade and Public Money}

Public-money language hides the chain.

Funds taken from Paul to pay Peter are not described as Paul's money. They are described as public revenue. Later, funds taken from Joseph to pay Paul are also described as public revenue. The continuity of extraction disappears because each taking is routed through the treasury.

The treasury makes the cascade look like administration rather than serial taking.

But public money is not ownerless money. Each round has a source. The source is the next payer in the chain.

\section{The Cascade and Damage Mitigation}

Recipients in a cascade may be mitigating damage.

A person forced to contribute for decades may accept benefits later to reduce the harm inflicted on him. That acceptance may be understandable. It may be practically necessary. It may not imply endorsement.

But damage mitigation is not entitlement against the next payer.

The recipient can truthfully say:

\begin{quote}
I am trying to recover from a system that harmed me.
\end{quote}

He cannot truthfully say, without more:

\begin{quote}
The next coerced payer owes me.
\end{quote}

The cascade turns damage mitigation into political entitlement. That conversion must be rejected.

\section{The Cascade and the Wrongdoer}

The wrongdoer remains central.

The state or coercive system creates the chain, enforces participation, defines eligibility, collects from current payers, and pays current beneficiaries. It then presents the transfer as if claimants and payers owe one another directly.

This misdescribes the relation.

Peter's claim is against the system that took from him. Paul's later claim is against the system that took from him. Joseph's later claim is against the system that took from him. The system avoids responsibility by making each cohort look to the next.

The wrongdoer cannot name the next payer.

\section{The Cascade Fallacy}

The cascade fallacy has this form:

\begin{enumerate}
  \item \(C_n\) was previously coerced to pay.
  \item Therefore \(C_n\) is owed.
  \item Therefore \(C_{n+1}\) may be coerced to pay \(C_n\).
  \item Because \(C_{n+1}\) paid, \(C_{n+1}\) is later owed.
  \item Therefore \(C_{n+2}\) may be coerced to pay \(C_{n+1}\).
\end{enumerate}

The error occurs at each transition from claimant to next payer.

The fact that \(C_n\) was wronged does not establish that \(C_{n+1}\) owes. Each step assumes what must be proven. The cascade survives by repeating the same invalid inference.

\section{Why the Cascade Feels Fair}

The cascade feels fair because each participant can tell a partial truth.

The beneficiary says:

\begin{quote}
I paid in.
\end{quote}

The payer is told:

\begin{quote}
You will receive later.
\end{quote}

The system says:

\begin{quote}
Everyone participates in turn.
\end{quote}

These truths are partial because they omit consent and source. Paying under compulsion is not voluntary contribution. Receiving later from new compelled payers is not clean repayment. Participation in turn is not consent if the rotation is enforced.

The cascade feels fair because it presents the sequence as reciprocity while hiding coercion at each link.

\section{Why the Cascade Is Politically Stable}

The cascade is politically stable because it creates defenders.

Each cohort that was harmed becomes invested in later extraction. The prior victim becomes current beneficiary. The beneficiary defends the system because his benefit depends on its continuation. The system thereby turns victims into advocates.

This is the genius of the cascade. It does not merely take. It converts the taken-from into future defenders of taking.

The prior victim's anger is redirected away from the system and toward preservation of the system, because only the system promises recovery.

\section{The Moral Cost of Cessation}

Cessation has moral cost.

It leaves some people unpaid. It exposes reliance. It disappoints expectations. It reveals that prior contributions were not stored for the contributors. It may produce hardship. It may feel like betrayal.

Those costs are real.

But continuation also has moral cost. It forces new people to pay. It creates the next grievance. It deepens reliance. It enlarges the chain. It pushes the final injury forward.

The choice is not between harm and no harm. It is between recognizing the current residual harm and manufacturing future residual harm.

\section{The Cleanest Possible Cessation}

The cleanest possible cessation follows several rules.

\begin{enumerate}
  \item Stop new compulsory entry where possible.
  \item Stop representing future theft as earned repayment.
  \item Identify the final aggrieved class honestly.
  \item Use clean sources, if any exist, for partial repair.
  \item Return identifiable funds where possible.
  \item Permit private alternatives immediately.
  \item Refuse to create new victims for the sake of completing an impossible chain.
\end{enumerate}

These measures do not make the system's history clean. They prevent additional propagation.

\section{Residual Wrong}

Cessation leaves residual wrong.

Residual wrong is the unpaid injury remaining when an immoral chain stops. It is not a refutation of cessation. It is evidence of the chain's prior immorality.

The residual wrong should be named:

\begin{quote}
This class was extracted from and will not receive the promised future extraction from others.
\end{quote}

That statement is hard. It is also more honest than continuing to pretend the system can make everyone whole by finding the next payer.

Residual wrong is tragedy. Continuation is multiplication.

\section{The No-New-Victim Rule}

The no-new-victim rule is the governing principle of cessation.

\begin{quote}
A prior uncompensated injury may not be used to justify creating a new uncompensated injury.
\end{quote}

This rule does not deny the prior injury. It denies the proposed method of addressing it.

If no clean source exists, then full restitution may be unavailable. That is a painful conclusion. But the absence of clean source does not authorize dirty source. The absence of clean repair does not authorize new theft.

\section{Objection: Stopping Cheats the Last Contributors}

Stopping does leave the last contributors with a grievance. But continuation cheats later contributors.

The objection identifies a real injury and draws the wrong conclusion. It says that because the final contributors are harmed, another class must be harmed to pay them. That repeats the cascade fallacy.

The honest answer is:

\begin{quote}
Yes, the last contributors were wronged. No, that wrong does not authorize wronging the next class.
\end{quote}

\section{Objection: The Chain Is Mutual Obligation}

A defender may say the chain is not theft but mutual obligation across generations.

This would require consent. Without voluntary entry and exit, the chain is not a covenant. It is compulsory rotation. Each cohort is told to pay because prior cohorts paid, and to expect later payment from future cohorts.

Mutual obligation cannot be manufactured by coercion. A compelled chain is not a contract.

\section{Objection: Everyone Benefits from Continuation}

Some benefit from continuation. Others are compelled to fund it. Future persons are burdened. The existence of benefit does not settle title.

The cascade's defenders point to beneficiaries and ignore the next coerced class. But every continuation has both sides. Benefit to one link is funded by injury to the next.

A chain cannot be judged only by the link receiving payment.

\section{Objection: Cessation Is Politically Impossible}

Political difficulty does not alter moral structure.

A system may be hard to stop because beneficiaries are numerous, expectations are entrenched, and the cost of honesty is high. That explains persistence. It does not justify continuation.

A wrong does not become right because many people depend on its continuation.

\section{Objection: The State Must Honor Its Promises}

The state should not make promises that require future coercion.

If the state promised Peter payment from Paul without Paul's consent, the promise itself was defective. A promisor cannot bind an unwilling stranger as future payer. The state's obligation to Peter does not become Paul's obligation merely because the state promised Peter that Paul would be made to pay.

A promise funded by future theft is not a clean promise.

\section{Objection: Stopping Destroys Trust}

Continuing a dirty chain preserves false trust.

It teaches people to trust that the next victim will be found. It treats moral legitimacy as a function of political continuity. It rewards the belief that if one is extracted from long enough, one earns the right to extract from others.

Cessation damages trust in the chain. That trust deserves to be damaged.

A system should not be trusted to pay old claims by creating new ones.

\section{Objection: Reform Can Fix the Chain}

Reform can reduce harm, but reform that preserves the next-victim structure does not fix the chain.

Changing rates, ages, formulas, eligibility, funding labels, or trust-fund accounting may alter distribution. It does not resolve the moral defect if current recipients are still funded by new coerced payers who do not owe them.

A cascade with better accounting is still a cascade.

\section{Relation to Prior Work}

This paper belongs to a sequence on extraction, receipt, restitution, liability, source, title, and consent.

Prior extraction-side analysis classifies taxation as coercive taking under legislative form. Receipt-side analysis argues that prior victimization does not sanitize receipt of currently extracted proceeds. Companion papers argue that pay-as-you-go transfers are corrupt at inception, restitution by new crime is not restitution, reparations funded by new extraction fail the source requirement, good distribution does not clean dirty funds, sympathetic recipients do not create title, state debt does not automatically become taxpayer debt, damage mitigation is not entitlement, the wrongdoer cannot name the payer, legal entitlement is not moral title, public money is not ownerless money, a claim is not a blank check, the beneficiary becomes the alibi, and dirty exit does not prove consent.

The present paper adds the cascade-side result:

\begin{quote}
Serial restitution by future theft cannot be morally completed.
\end{quote}

It explains why stopping the chain, despite leaving a final residual injury, is cleaner than continuing the chain through the next victim.

\section{Conclusion}

The entitlement cascade is the moral structure of serial coercive repair.

One class is wronged. The next class is coerced to compensate it. That next class is then wronged. A later class is coerced to compensate it. Each payment appears to repair the prior injury, but each payment creates the next injury. The chain never reaches clean restitution because it never uses clean source.

The hardest fact is the final aggrieved class. If the chain stops, someone is left standing. That person or class may have paid in, relied, and expected future payment. The grievance is real.

But that grievance is not a warrant to create another one.

Continuing the chain does not repair the injustice. It propagates it. It hides the final victim by moving him into the future. It turns every prior extraction into an argument for the next extraction. It keeps the music playing by guaranteeing that someone else will be left without a chair later.

The clean answer is not painless. It is cessation.

Stop the chain. Name the residual wrong. Use clean sources if any exist. Refuse to create new victims to disguise old injury.

The final rule is plain:

\begin{quote}
A system that requires a next victim to compensate the prior victim has no clean moral completion by continuation.
\end{quote}

\begin{thebibliography}{9}

\bibitem{mckinley-taxation}
J.\ C.\ McKinley,
\textit{Taxation as Legislated Crime via Majority Hegemony},
Zenodo (2026).
\href{https://doi.org/10.5281/zenodo.18530711}{https://doi.org/10.5281/zenodo.18530711}.

\bibitem{mckinley-receipt}
J.\ C.\ McKinley,
\textit{Receipt of Legislated Theft: Why Prior Victimization Does Not Sanitize Currently-Extracted Proceeds},
manuscript (2026).

\bibitem{mckinley-firstbeneficiaries}
J.\ C.\ McKinley,
\textit{The First Beneficiaries: Why Pay-As-You-Go Transfers Are Corrupt at Inception},
manuscript (2026).

\bibitem{mckinley-newcrime}
J.\ C.\ McKinley,
\textit{Restitution by New Crime Is Not Restitution: A General No-Go on Repairing Prior Wrong Through Fresh Coercion},
manuscript (2026).

\bibitem{mckinley-reparations}
J.\ C.\ McKinley,
\textit{Reparations from Stolen Funds: Why Historical Wrong Does Not Authorize Current Extraction},
manuscript (2026).

\bibitem{mckinley-aftertheft}
J.\ C.\ McKinley,
\textit{After the Theft: Why Good Distribution Does Not Clean Dirty Funds},
manuscript (2026).

\bibitem{mckinley-sympathy}
J.\ C.\ McKinley,
\textit{Sympathetic Recipients and the Laundering of Coercion: Why Need, Victimhood, and Moral Attractiveness Do Not Create Title},
manuscript (2026).

\bibitem{mckinley-stateowes}
J.\ C.\ McKinley,
\textit{The State Owes You; I Do Not: A No-Go on Converting State Debt into Taxpayer Debt},
manuscript (2026).

\bibitem{mckinley-damage}
J.\ C.\ McKinley,
\textit{Damage Mitigation Is Not Entitlement: A No-Go on Recasting Survival Inside Coercion as Clean Moral Title},
manuscript (2026).

\bibitem{mckinley-wrongdoer}
J.\ C.\ McKinley,
\textit{The Wrongdoer Cannot Name the Payer: A No-Go on Assigning Liability to Uninvolved Third Parties},
manuscript (2026).

\bibitem{mckinley-legalentitlement}
J.\ C.\ McKinley,
\textit{Legal Entitlement Is Not Moral Title: A No-Go on Treating Statutory Eligibility as Ownership},
manuscript (2026).

\bibitem{mckinley-publicmoney}
J.\ C.\ McKinley,
\textit{Public Money Is Not Ownerless Money: A No-Go on Linguistic Laundering Through the Treasury},
manuscript (2026).

\bibitem{mckinley-claim}
J.\ C.\ McKinley,
\textit{A Claim Is Not a Blank Check: Why Injury, Reliance, or Need Does Not Authorize Any Available Payer},
manuscript (2026).

\bibitem{mckinley-alibi}
J.\ C.\ McKinley,
\textit{The Beneficiary Becomes the Alibi: A No-Go on Using Visible Help to Excuse Invisible Coercion},
manuscript (2026).

\bibitem{mckinley-cleanexit}
J.\ C.\ McKinley,
\textit{The Clean Exit Problem: How Coercive Systems Remove Alternatives and Then Call Participation Consent},
manuscript (2026).

\bibitem{helvering}
\textit{Helvering v.\ Davis},
301 U.S. 619 (1937).

\bibitem{flemming}
\textit{Flemming v.\ Nestor},
363 U.S. 603 (1960).

\bibitem{madoff}
\textit{Picard v.\ Katz} and related proceedings,
\textit{In re Bernard L. Madoff Investment Securities LLC},
U.S. Bankruptcy Court for the Southern District of New York.

\bibitem{nozick}
R.\ Nozick,
\textit{Anarchy, State, and Utopia},
Basic Books (1974).

\bibitem{locke}
J.\ Locke,
\textit{Second Treatise of Government},
Awnsham Churchill, London (1689).

\bibitem{barnett}
R.\ E.\ Barnett,
\textit{The Structure of Liberty: Justice and the Rule of Law},
Oxford University Press (1998).

\bibitem{epstein}
R.\ A.\ Epstein,
\textit{Takings: Private Property and the Power of Eminent Domain},
Harvard University Press (1985).

\end{thebibliography}

\end{document}